\section{Introduction}
\label{sec:intro}

Practitioners size serving hardware by parameter count: count parameters,
multiply by bytes per weight, buy that much memory. This habit breaks down for
long-context and multi-sequence serving. There, the
binding constraint is the key--value (KV) cache that attention layers append to
at every decoded token. Serving systems already treat it as the scarce
resource: PagedAttention exists because per-request KV memory dominates and
fragments GPU memory \cite{kwon2023paged}, DistServe replicates the KV cache
when it splits prefill from decoding \cite{zhong2024distserve}, and
FlashAttention restructures the kernel around memory traffic \cite{dao2022flashattention}.

The KV cache does not scale with parameter count. Its size follows three
architectural choices. \emph{Attention geometry}: a full-attention layer stores
keys and values for $2 \times n_{\mathrm{kv\,heads}} \times d_{\mathrm{head}}$
elements per token, so layers, KV heads, and head dimension set the marginal
bytes. \emph{Attention scope}: a sliding-window layer retains only its window,
not the whole context \cite{jiang2023mistral}. \emph{State type}: recurrent and
linear-attention layers carry a fixed-size state instead of a growing cache
\cite{gu2023mamba,katharopoulos2020linear}, and multi-head latent attention
compresses keys and values into a low-rank latent \cite{deepseekai2024v2}.
Grouped-query and multi-query attention cut KV heads directly
\cite{shazeer2019mqa,ainslie2023gqa}, and cross-layer sharing stores one cache
for many layers \cite{sun2024yoco}. Two checkpoints with equal parameter counts
can therefore differ by an order of magnitude in KV bytes per token. On
unified-memory edge devices, where weights and KV compete for one pool of
memory, this difference decides how many concurrent long contexts fit; measured
runtime comparisons on such hardware already rank models by KV and prompt-cache
behavior \cite{rajesh2025edgemem}.

Which of these designs the field actually ships is an empirical question, and
the answer changes what hardware practitioners should buy. Answering it needs a
census: open weights so the config is public, published configs so geometry is
exact, and release dates so trends are datable. We built one from the Hugging
Face Hub, restricted to original publishers and decoder-only language models
released since January 2023. The census covers \kvcCensusRows{} checkpoints from
\kvcOrgCount{} publishers; after deduplicating base/instruct twins that share
geometry, \kvcDedupRows{} remain. We skipped \kvcGatedSkipped{} gated
repositories whose configs are not publicly downloadable and report this as a
coverage limitation. The deduplicated set spans full attention (MHA, MQA, GQA),
sliding-window and local/global mixes, latent attention, recurrent hybrids
(Mamba, RWKV, gated DeltaNet), and cross-layer KV sharing. Table~\ref{tab:overview}
summarizes the census by year.

We ask four research questions and report four pre-registered verdicts
(thresholds fixed before trend inspection):
\begin{itemize}
  \item \textbf{RQ1 (trend).} How did KV bytes per token, and KV bytes per token
  per billion parameters, evolve from 2023 to 2026? \emph{H1: the median
  per-parameter KV cost fell substantially} --- verdict \kvcHOneVerdict{}, the
  pooled 2025--2026 median is \kvcHOneRatio{} of the 2023 median.
  \item \textbf{RQ2 (mechanisms).} How much cross-model variance do GQA ratio,
  head dimension, depth, window share, recurrent share, latent attention, and
  cross-layer sharing explain, relative to parameter count? \emph{H2: parameter
  count is a weak predictor} --- verdict \kvcHTwoVerdict{}, $R^2=\kvcParamsRsq$
  for parameters alone versus $R^2=\kvcGeomRsq$ for geometry and mechanisms
  alone. \emph{H3: mixture-of-experts (MoE) models have much lower KV per total
  parameter} --- verdict \kvcHThreeVerdict{}, the MoE median is
  \kvcHThreeRatio{} of the dense median.
  \item \textbf{RQ3 (capacity).} Under fixed memory budgets and weight
  precisions, what context and concurrency can each model serve? We define the
  \emph{KV-parity context} $T^{*}$: the context length at which one sequence's
  KV bytes equal the model's weight bytes.
  \emph{H4: the KV-parity context grew by orders of magnitude} --- verdict
  \kvcHFourVerdict{}, the median grew only \kvcHFourRatio$\times$.
  \item \textbf{RQ4 (validity).} Do analytic estimates match the cache tensors
  a reference implementation (HF \texttt{transformers}) actually allocates?
  Validation covers \kvcValRows{} models in \kvcValClasses{} attention classes;
  \kvcValWithin{} rows agree within 1\%, and we explain every larger gap.
\end{itemize}

Our contributions:
\begin{itemize}
  \item A census of \kvcCensusRows{} open-weight checkpoints with per-checkpoint
  KV accounting that covers every attention class we found on the Hub, released
  as data, code, and reproducible pipeline.
  \item Population-level trend and regression results: per-parameter KV cost is
  falling, but it is set by architecture, not size; MoE amplifies this
  separation between parameters and memory.
  \item A capacity model that turns each config into context and concurrency
  limits under stated budgets and precisions, with the KV-parity context as a
  single comparable number per model.
  \item A validation experiment checking analytic bytes against allocated cache
  tensors, including the cases where implementations diverge from the
  theoretical minimum.
  \item An honest negative result: H4 is refuted. The median KV-parity context
  grew only \kvcHFourRatio$\times$, not by orders of magnitude. We report the
  refutation as a first-class finding, not a failure.
\end{itemize}

The one-sentence thesis is \emph{memory is architecture}: to know what a
model costs to serve, read its config, not its parameter count.
